\section{Discussions}
\label{sec:discussions}

In this work, we introduced \method, an online-MARL-inspired alignment method that encourages diverse generation while preserving response quality. \method~serves as a drop-in replacement for existing post-training alignment pipelines, requiring no architectural modifications to the underlying model. Compared to previous methods, its dual-mode design enables seamless switching between producing a single high-confidence answer and generating a diverse set of plausible responses, offering greater flexibility depending on the downstream needs. The ability to surface conceptually distinct yet high-quality responses has broad implications for high-stakes domains such as medical diagnosis, legal reasoning, and scientific ideation.

\textbf{Limitations and Future Work.} Several limitations point to promising directions for future research. First, our quality reward model (SkyReward-V2) tends to favor verbose responses, even for prompts that should be answered briefly, and our uniform length penalty insufficiently addresses this across prompts of varying complexity. For instance, when training GLM-4-9B with \method, we observed suboptimal native quality score on NoveltyBench compared to the base model because its reward model doesn't penalize verbosity therefore it reward hack by adding unnecessary preambles such as ``As a model.../To answer these questions as a model...'' Removing these preambles  recovers the performance from 2.739 to 3.302 (+20\%). Future work should explore adaptive length penalties or alternative reward models that better reflect conciseness. Second, because \method~introduces no external data during training, generative diversity is bounded by the initial policy's capacity, and incorporating retrieval augmentation or external data sources could meaningfully expand the exploration space. Finally, there is a risk of producing more varied but misleading, unsafe, or unsupported outputs, especially in high-stakes domains such as medical diagnosis and legal reasoning. Although \method~mitigates this through its dual-mode design, whose standard mode retains the model's general capability, future work can further combine the diverse mode with verification, calibration, retrieval, or domain-specific safeguards.